%=======================================================================
\section*{Main [Yifu]}
%=======================================================================
Air connectivity is an explicit object of national policy. Governments
subsidise routes, expand airport capacity, and negotiate air-service
agreements on the premise that connections to the global air network raise
trade, tourism, and productivity, and a large literature documents such
gains [TODO cite]. Aviation is also among the most difficult transport
sectors to decarbonise: it has no near-term electrification path, and its
emissions are governed by separate international accounting conventions
rather than by national inventories [TODO cite]. The carbon consequences of
connectivity policy are therefore a first-order question, and one for
which, to our knowledge, no causal estimate exists at the global level.

The direction of the effect is not in serious doubt: more connections mean
more flying. The empirical questions concern magnitude and structure.
First, proportionality: does a one percent gain in connectivity raise
emissions by more or less than one percent, and through which margin of
aviation activity? Network development changes not only the amount of
flying but its organisation, since hub formation can consolidate traffic
into fewer, fuller, and more efficient movements. Second, incidence: which
countries and which airports carry the response, and does it spill across
borders through the network? Third, accounting: the emissions generated by
connectivity growth need not be booked where they are generated, because
international aviation is attributed under bunker conventions rather than
territorially [TODO cite].

We answer these questions with two inputs (Methods). The first is a new
emissions panel: flight-stage-resolved CO$_2$ computed from OAG schedule
data under the EEA/EMEP Guidebook methodology for 6{,}356 airports over
1996--2023, disaggregated by flight stage and by domestic and international
traffic, so that national emissions can be constructed under any allocation
rule without double counting and so that the response can be decomposed on
the identity used to build the data. The second is identification: we
instrument a country's connectivity, measured by the Global Air
Connectivity Index (GACI), with the interaction of world aviation
technology and the country's air-versus-sea market-access geography in
1996, in the spirit of Feyrer [TODO cite]. The instrument is strong
(first-stage Kleibergen--Paap $F$ of 18 with standard errors clustered by
country, 154 with heteroskedasticity-robust errors), survives a stress test that
interacts the global cycle with baseline population, income, and airport
capacity, and passes three exclusion-restriction tests: non-aviation
emissions do not respond to it, air geography rather than sea geography
carries its effect, and where it does not move connectivity it does not
move emissions (Methods). The analysis is
organised in five layers: the causal effect at the country and airport
levels, its decomposition on the accounting identity, its heterogeneity,
its cross-border spillovers, and its policy accounting.

